\chapter{Spektrum von Signalfolgen}\label{sec:spektrum-von-signalfolgen}

Die Fourier-Transformation erlaubt es, zeitkontinuierlichen Signalen ein Spektrum und zeitkontinuierlichen Systemen einen Frequenzgang zuzuordnen. Die Fourier-Transformation von Signalfolgen ist für zeitdiskrete Anwendungen das Äquivalent zur bereits diskutierten Fourier-Transformation für zeitkontinuierliche Signale. Sie wird in diesem Kapitel eingeführt und ist wesentliche Voraussetzung für die Beschreibung zeitdiskreter Signale sowie die Analyse zeitdiskreter Systeme im Frequenzbereich.

Die Fourier-Transformierte von Signalfolgen kann über ihre Definitionsgleichung bestimmt werden. Wie bei der z-Transformation kann die eher aufwendige Bestimmung von Korrespondenzen über die Definitionsgleichung vermieden werden, wenn die vorliegende Signalfolge mit Rechenregeln auf Folgen mit bekannten Korrespondenzen zurückgeführt wird. Deshalb werden für die Fourier-Transformation von Signalfolgen die gebräuchlichen Rechenregeln hergeleitet und der Nutzen an Beispielen aufgezeigt. Weiterhin weist die Fourier-Transformation von Signalfolgen Symmetrien auf, die berechnet und an einem Beispiel verdeutlicht werden.

Abschließend wird ein Zusammenhang zwischen der Fourier-Transformation von Signalfolgen, der zeitkontinuierlichen Fourier-Transformation und der z-Transformation hergestellt.

\section{Grundlagen}\label{sec:7.1}

\subsection{Eigenfunktionen zeitdiskreter LTI-Systeme}\label{sec:eigenfunktionen-zeitdiskreter-lti-systeme}

Lineare Differenzengleichungen mit konstanten Koeffizienten beschreiben zeitdiskrete LTI-Systeme. Die Lösungen dieser Differenzengleichungen setzen sich aus Exponentialfolgen zusammen. Eine besondere Stellung nehmen komplexe Exponentialfolgen der Form
\begin{align}
x[ k] &= e^{j  \Omega_{0}  k} \label{eq:7.1}
\end{align}
ein. Sie werden als Eigenfolgen des Systems bezeichnet. Bei einer Anregung mit diesen Folgen reagiert das System mit einer Folge gleicher normierter Frequenz \(Ω_{0}\) im Allgemeinen aber unterschiedlicher Amplitude und Phase.

\begin{example}{Anregung eines Systems 1. Ordnung mit einer komplexen Exponentialfolge}
\label{beispiel-anregung-eines-systems-1.-ordnung-mit-einer-komplexen-exponentialfolge}

Ein System mit der Übertragungsfunktion
\begin{align}
G(z) &= \frac{z}{z - \lambda} \label{eq:7.2}
\end{align}
wird mit einem Signal
\begin{align}
x[ k] &= e^{j  \Omega_{0}  k}  \sigma[ k] \label{eq:7.3}
\end{align}
angeregt. Das Eingangssignal hat die z-Transformierte
\begin{align}
X(z) &= \frac{z}{z - e^{j  \Omega_{0}}} \label{eq:7.4}
\end{align}
Die z-Transformierte des Ausgangssignals lautet nach den Darstellungen zur z-Transformation
\begin{align}
Y(z) &= G(z)  X(z) = \frac{z}{z - \lambda}  \frac{z}{z - e^{j  \Omega_{0}}} \label{eq:7.5}
\end{align}
Über Partialbruchzerlegung und Rücktransformation ergibt sich die Ausgangsfolge \(y[k]\) mit
\begin{align}
{y\left[k\right]} &=  {\frac{\lambda ^{k+1} }{\lambda -e^{j \Omega _{0} } }  \sigma \left[k\right]+\frac{e^{j \Omega _{0} } }{e^{j \Omega _{0} } -\lambda }  e^{j \Omega _{0}  k}  \sigma \left[k\right]=\frac{\lambda ^{k+1} }{\lambda -e^{j \Omega _{0} } }  \sigma \left[k\right]+\left. G\left(z\right)\right|_{z=e^{j \Omega _{0} } }  e^{j \Omega _{0}  k}  \sigma \left[k\right]} \\
{} &=  {\frac{\lambda ^{k+1} }{\lambda -e^{j \Omega _{0} } }  \sigma \left[k\right]+A\left(\Omega _{0} \right) e^{j \varphi \left(\Omega _{0} \right)}  e^{j \Omega _{0}  k}  \sigma \left[k\right]}
 \label{eq:7.6}
\end{align}
Das System reagiert auf die harmonische Anregung mit der Kreisfrequenz \(Ω_{0}\) abgesehen von Einschwingvorgängen mit einem harmonischen Signal derselben Kreisfrequenz. Für den Fall einer harmonischen Anregung müssen demnach nur das Verhältnis der Ein- und Ausgangsamplitude sowie die Phasenverschiebung bestimmt werden. Beide Größen sind von der normierten Kreisfrequenz \(Ω_{0}\) abhängig. Wie bei zeitkontinuierlichen Signalen wird das Verhältnis der Amplituden als Amplitudengang \(A(Ω)\) und die Phasenverschiebung als Phasengang ϕ(Ω) bezeichnet.

\end{example}

Das Ausgangssignal ergibt sich aus dem Produkt des Eingangssignals mit einem komplexen Faktor, der zu einer Änderung der Amplitude und der Phase führt. Die Amplitude ergibt sich aus dem Betrag der Übertragungsfunktion \(G(z)\) an der Stelle z = \(e^{jΩo}\). Die Phase ergibt sich aus der Phase der Übertragungsfunktion \(G(z)\) an der Stelle z = \(e^{jΩo}\). Die Systemreaktion kann damit bei Anregung eines Systems mit einer komplexen Exponentialfunktion vergleichsweise einfach bestimmt werden.

Auch nicht periodische Signale lassen sich als Linearkombination komplexer Exponentialfolgen beschreiben. Für jede einzelne Folge kann die Wirkung des Systems berechnet werden. Das Ausgangssignal ergibt sich aus der Linearkombination von komplexen Exponentialfolgen mit geänderten Amplituden und geänderten Phasen. Die Darstellung von Signalfolgen als Linearkombination komplexer Exponentialfolgen wird als Fourier-Transformation für Signalfolgen bezeichnet. Sie ist Gegenstand dieses Kapitels.

\subsection{Definitionsgleichung der Fourier-Transformation von Signalfolgen}\label{sec:7.1.2}

Das Spektrum eines zeitkontinuierlichen Signals wird nach den Ausführungen in Teil A des Skriptes berechnet über das Integral
\begin{align}
X(\omega) &= \int_{- \infty}^{\infty}{x(t)  \ }e^{- j  \omega  t}\ dt \label{eq:7.7}
\end{align}
Ein abgetastetes Signal \(x_{A}(t)\)
\begin{align}
x_{A}(t) &= \sum\limits_{k = - \infty}^{\infty}{x(t)  \delta\left( t - k  T_{A} \right)} = \sum\limits_{k = - \infty}^{\infty}{x[ k]  \delta\left( t - k  T_{A} \right)} \label{eq:7.8}
\end{align}
ist zu jedem Zeitpunkt definiert. Sein Spektrum berechnet sich mit Gleichung \eqref{eq:7.7} zu
\begin{align}
X_{A}(\omega) &= \int_{- \infty}^{\infty}{x_{A}(t)  \ }e^{- j  \omega  t}\ dt = \int_{- \infty}^{\infty}{\left( \sum\limits_{k = - \infty}^{\infty}{x(t)  \delta\left( t - k  T_{A} \right)} \right)  \ }e^{- j  \omega  t}\ dt \label{eq:7.9}
\end{align}
Die Reihenfolge von Summation und Integration kann vertauscht werden. Mit der Ausblendeigenschaft der Impulsfunktion kann der Ausdruck umgeformt werden zu
\begin{align}
{X_{A} \left(\omega \right)} &=  {\int _{-\infty }^{\infty }\left(\sum\limits _{k=-\infty }^{\infty }x\left[k\right] \delta \left(t-k T_{A} \right) \right) {\; } e^{-j \omega  t} {\; }dt=\sum\limits _{k=-\infty }^{\infty }x\left[k\right] \int _{-\infty }^{\infty }\delta \left(t-k T_{A} \right) {\; } e^{-j \omega  t} {\; }dt } \\
{} &=  {\sum\limits _{k=-\infty }^{\infty }x\left[k\right] e^{-j \omega  k T_{A} }  {\; } \int _{-\infty }^{\infty }\delta \left(t-k T_{A} \right){\; } dt=\sum\limits _{k=-\infty }^{\infty }x\left[k\right] e^{-j \omega  k T_{A} }  }
 \label{eq:7.10}
\end{align}
Gleichung \eqref{eq:7.10} weist der Folge von Abtastwerten \(x[k]\) ein Spektrum \(X_{A}(ω)\) zu. Dabei tritt die Frequenz ω immer in Kombination mit der Abtastzeit \(T_{A}\) auf. Deshalb wird bei Signalfolgen das Spektrum als Funktion der normierten Kreisfrequenz
\begin{align}
\Omega &= \omega  T_{A} = \frac{\omega}{f_{A}} \label{eq:7.11}
\end{align}
dargestellt. Wegen der Periodizität der Exponentialfunktion mit imaginärem Argument ist jeder Summand der Summe in Gleichung \eqref{eq:7.10} periodisch in 2∙π. Deshalb ist auch das Spektrum der Folge \(x[k]\) periodisch in 2∙π. Es ist damit ausreichend, das Spektrum in den Grenzen von - π \ldots{} + π darzustellen. Durch die Normierung über die Abtastzeit \(T_{A}\)
\begin{align}
\omega &= \frac{\Omega}{T_{A}} \label{eq:7.12}
\end{align}
ergibt sich die in Tabelle \ref{tab:7.1} gezeigte Zuordnung zwischen der normierten Frequenz Ω und der Kreisfrequenz ω.

\begin{table}
  \centering
  \caption{Zuordnung zwischen der normierten Frequenz Ω und der Kreisfrequenz ω bei der Fourier-Transformation für Folgen}
  \label{tab:7.1}
  \begin{tabular}{ll}
    \toprule
    Normierte Kreisfrequenz Ω & Kreisfrequenz ω \\
    \midrule
    - π & - \(ω_{A}\)/2 \\
    - π/2 & - \(ω_{A}\)/4 \\
    0 & 0 \\
    π/2 & \(ω_{A}\)/4 \\
    π & \(ω_{A}\)/2 \\
    \bottomrule
  \end{tabular}
\end{table}

Damit ergibt sich die Definitionsgleichung der Fourier-Transformation von Signalfolgen zu
\begin{align}
X(\Omega) &= \sum\limits_{k = - \infty}^{\infty}{x[ k]  e^{- j  \Omega  k}} \label{eq:7.13}
\end{align}
Die Fourier-Transformation für Signalfolgen transformiert beliebige Signalfolgen \(x[k]\) in eine kontinuierliche, komplexe Funktion der reellen Variable Ω. Wie bei zeitkontinuierlichen Signalen wird die Funktion \(X(Ω)\) als Spektrum der Signalfolge \(x[k]\) bezeichnet. Auch die Schreibweise
\begin{align}
F\left\{ x[ k] \right\} &= X(\Omega) \label{eq:7.14}
\end{align}
sowie das Hantel-Symbol
\begin{align}
x[ k] \origof X(\Omega) \label{eq:7.15}
\end{align}
werden von der Fourier-Transformation für zeitkontinuierliche Signale übernommen. Die Fourier-Transformierte kann in kartesischen oder Polar-Koordinaten dargestellt werden.
\begin{align}
X(\Omega) &= X_{R}(\Omega) + j  X_{I}(\Omega) = \left| X(\Omega) \right|  e^{j  \varphi(\Omega)} \label{eq:7.16}
\end{align}
Bei der Darstellung in kartesischen Koordinaten wird die Fourier-Transformierte in Realteil \(X_{R}(Ω)\) und Imaginärteil \(X_{I}(Ω)\) zerlegt. Bei der Darstellung in Polarkoordinaten wird die komplexe Fourier-Transformierte mit Betrag \textbar \(X(Ω)\)\textbar{} und Phase ϕ(Ω) dargestellt.

\subsection{\texorpdfstring{Fourier-Transformation von grundlegenden Signalfolgen }{Fourier-Transformation von grundlegenden Signalfolgen }}\label{fourier-transformation-von-grundlegenden-signalfolgen}

Wie bei der z-Transformation werden zunächst einige Korrespondenzen der Fourier-Transformation von Signalfolgen direkt über die Definitionsgleichung \eqref{eq:7.13} berechnet. Dabei ergibt sich das Spektrum einer Signalfolge durch Auswertung einer in der Regel unendlichen, oft aber endlichen Reihe.

\subsubsection{Impulsfolge}\label{sec:impulsfolge-1}

Aus der Definition der diskreten Impulsfolge
\begin{align}
x[ k] = \delta[ k] = \left\{ \begin{aligned}
 & 1\text{  f}\text{ü}\text{r k} = 0 \\
 & \text{0 f}\text{ü}\text{r ganzzahlige  k} \neq 0
\end{aligned} \right.\ \label{eq:7.17}
\end{align}
ergibt sich das Spektrum durch Einsetzen in die Definitionsgleichung zu
\begin{align}
X(\Omega) &= \sum\limits_{k = - \infty}^{\infty}{\delta[ k]  e^{- j  \Omega  k}} = e^{- j  0  \Omega}  \sum\limits_{k = - \infty}^{\infty}{\delta[ k]} = 1 \label{eq:7.18}
\end{align}
Bild \ref{fig:7.1} stellt die Impulsfolge und ihr Spektrum gegenüber. Wie bereits bei der zeitkontinuierlichen Fourier-Transformation werden zur Darstellung eines Impulses die Frequenzen im gesamten Spektralbereich benötigt.

\begin{figure}
  \centering
  \includegraphics[alt={Impulsfolge und ihr Spektrum}]{new media/7_1_impulsfolge_und_ihr_spektrum.png}
  \caption{Impulsfolge und ihr Spektrum}
  \label{fig:7.1}
\end{figure}

Wird der Impuls um \(k_{0}\) verschoben, ändert sich die Fourier-Transformierte zu
\begin{align}
X(\Omega) &= \sum\limits_{k = - \infty}^{\infty}{\delta\left[ k - k_{0} \right]  e^{- j  \Omega  k}} = e^{- j  \Omega  k_{0}}  \sum\limits_{k = - \infty}^{\infty}{\delta\left[ k - k_{0} \right]} = e^{- j  \Omega  k_{0}} \label{eq:7.19}
\end{align}
Eine Verschiebung des Impulses um \(k_{0}\) \textgreater{} 0 nach rechts führt zu einer Multiplikation der Fourier-Transformierten mit z = \(e^{-jΩko}\). Diese Multiplikation ändert die Phase des Spektrums, sein Betrag bleibt dagegen konstant.

\subsubsection{Kausale Rechteckfolge}\label{sec:kausale-rechteckfolge}

Für die Rechteckfolge mit der Definition
\begin{align}
x[ k] &= \sigma[ k] - \sigma[ k - K] \label{eq:7.20}
\end{align}
wird die Summe in der Definitionsgleichung der Fourier-Transformation von Signalfolgen endlich. Mit der endlichen geometrischen Reihe
\begin{align}
\sum\limits_{k = 0}^{K - 1}q^{k} &= \frac{1 - q^{K}}{1 - q} \label{eq:7.21}
\end{align}
ergibt sich die Fourier-Transformierte
\begin{align}
X(\Omega) &= \sum\limits_{k = 0}^{\infty}{\left( \sigma[ k] - \sigma[ k - K] \right)  e^{- j  \Omega  k}} = \sum\limits_{k = 0}^{K - 1}e^{- j  \Omega  k} = \frac{1 - e^{- j  \Omega  K}}{1 - e^{- j  \Omega}} \label{eq:7.22}
\end{align}
Zur übersichtlicheren Trennung von Betrag und Phase werden aus Nenner und Zähler Exponentialfunktionen ausgeklammert, sodass konjugiert komplexe Exponentialfunktionen entstehen. Sie werden zu einer Sinusfunktion zusammengefasst.
\begin{align}
X(\Omega) &= \frac{1 - e^{- j  \Omega  K}}{1 - e^{- j  \Omega}} = \frac{e^{- j  \frac{\Omega  K}{2}}}{e^{- j  \frac{\Omega}{2}}}  \frac{\sin\left( \frac{K  \Omega}{2} \right)}{\sin\left( \frac{\Omega}{2} \right)} = e^{- j  \frac{\Omega  (K - 1)}{2}}  \frac{\sin\left( \frac{K  \Omega}{2} \right)}{\sin\left( \frac{\Omega}{2} \right)} \label{eq:7.23}
\end{align}
Bild \ref{fig:7.2} stellt die Rechteckfolge für \(K_{1} = 5\) und \(K_{2} = 11\) und den zugehörigen Betrag des Spektrums gegenüber.

\begin{figure}
  \centering
  \includegraphics[alt={Rechteckfolge und Betrag ihres Spektrums für \(K_{1} = 5\) und \(K_{2} = 11\)}]{new media/7_2_rechteckfolge_und_betrag_ihres_s.png}
  \caption{Rechteckfolge und Betrag ihres Spektrums für \(K_{1} = 5\) und \(K_{2} = 11\)}
  \label{fig:7.2}
\end{figure}

Die erste Nullstelle \(Ω_{0}\) des Spektrums liegt an der Stelle
\begin{align}
\Omega_{0} &= \frac{2  \pi}{K} \label{eq:7.24}
\end{align}
Wird diese Nullstelle als Maß für die Breite des Spektrums angesehen, wird klar, dass das Spektrum mit sinkender Zahl K von Abtastwerten breiter wird. Im Extremfall liegt mit \(x[k]\) = \(δ[k]\) ein einziger Abtastwert vor, und das Spektrum wird mit \(X(Ω)\) = 1 beliebig breit.

\subsubsection{Kausale Exponentialfolge}\label{sec:kausale-exponentialfolge}

Die kausale Potenzfolge ist definiert als
\begin{align}
x[ k] &= \lambda^{k}  \sigma[ k] \label{eq:7.25}
\end{align}
Ihre Fourier-Transformierte ergibt sich aus der Definitionsgleichung zu
\begin{align}
X(\Omega) &= \sum\limits_{k = - \infty}^{\infty}{\lambda^{k}  \sigma[ k]  e^{- j  k  \Omega}} = \sum\limits_{k = 0}^{\infty}\left( \lambda  e^{- j  \Omega} \right)^{k} \label{eq:7.26}
\end{align}
Mit der Summationsformel für die geometrische Reihe
\begin{align}
\sum\limits_{k = 0}^{\infty}q^{k} &= \frac{1}{1 - q}\text{ für }|q|\  < \text{ 1} \label{eq:7.27}
\end{align}
ergibt sich für \textbar λ\textbar{} \textless{} 1 die Fourier-Transformierte
\begin{align}
F\left\{ \lambda^{k}  \sigma[ k] \right\} &= \sum\limits_{k = 0}^{\infty}\left( \lambda  e^{- j  \Omega} \right)^{k} = \frac{1}{1 - \lambda  e^{- j  \Omega}}\ \label{eq:7.28}
\end{align}
Die Fourier-Transformierte der kausalen Exponentialfolge kann auf diese Art nur für \textbar λ\textbar{} \textless{} 1 berechnet werden. Für \textbar λ\textbar{} ≥ 1 existiert sie nicht. Bild \ref{fig:7.3} stellt die kausale Exponentialfolgen mit \(λ = 0.9\) und λ = - 0.9 sowie den Betrag ihrer Spektren dar.

Sowohl für \(λ = 0.9\), als auch für λ = - 0.9 klingen die Signalfolgen \(x[k]\) ab. Wegen des negativen Vorzeichens handelt es sich bei λ = - 0.9 um eine alternierende Folge. Durch den schnellen Signalwechsel bei der alternierenden Signalfolge besitzt das Spektrum der Exponentialfolge mit λ = - 0.9 hohe Signalanteile bei den normierten Frequenzen Ω = ± π, die der Abtastfrequenz ± \(ω_{A}\)/2 entsprechen. Im Gegensatz dazu weist das Spektrum der Exponentialfolge mit \(λ = 0.9\) wesentliche Spektralanteile an der Stelle \(Ω = 0\) auf.

\begin{figure}
  \centering
  \includegraphics[alt={Kausale Exponentialfolge und Betrag ihres Spektrums für \(λ = 0.9\) und λ = - 0.9}]{new media/7_3_kausale_exponentialfolge_und_bet.png}
  \caption{Kausale Exponentialfolge und Betrag ihres Spektrums für \(λ = 0.9\) und λ = - 0.9}
  \label{fig:7.3}
\end{figure}

\subsection{Existenz der Fourier-Transformation von Signalfolgen}\label{sec:7.1.4}

Bei der Berechnung der Fourier-Transformierten der kausalen Exponentialfolge ergibt sich die Bedingung \textbar λ\textbar{} \textless{} 1 für die Konvergenz der zu berechnenden Reihe. Daraus resultiert die Frage, unter welchen Bedingungen die Fourier-Transformation von Signalfolgen konvergiert. Bei der Berechnung über die Definitionsgleichung wird die Reihe
\begin{align}
X(\Omega) &= \sum\limits_{k = - \infty}^{\infty}{x[ k]  e^{- j  \Omega  k}} \label{eq:7.29}
\end{align}
ausgewertet. Die Fourier-Transformierte existiert, wenn diese Reihe konvergiert. Eine hinreichende Bedingung für die Existenz der Fourier-Transformierten einer Folge ist ihre absolute Summierbarkeit.
\begin{align}
\sum\limits_{k = - \infty}^{\infty}{\left| x[ k] \right| < \infty} \label{eq:7.30}
\end{align}
Zwei Spezialfälle werden diskutiert, kausale Signalfolgen und zeitlich begrenzte Signalfolgen.

\subsubsection{Kausale Signalfolgen}\label{sec:kausale-signalfolgen}

Bei kausalen Signalfolgen kann die Konvergenzbedingung aus der geometrischen Reihe abgeleitet werden. Dabei wird ihre Konvergenz auf eine Betrachtung des Betrages von q reduziert. Um diesen Ansatz anwenden zu können, ist es wie bei der z-Transformation erforderlich, den Betrag der Folge \(x[k]\) mit Hilfe einer Exponentialfolge nach oben abzuschätzen.
\begin{align}
\left| x[ k] \right| \leq X_{MAX}  r^{k} \label{eq:7.31}
\end{align}
Mit der Abschätzung der Folge \(x[k]\) gilt für \(X(Ω)\)
\begin{align}
\left| X(\Omega) \right| &= \left| \sum\limits_{k = 0}^{\infty}{x[ k]  e^{- j  \Omega  k}} \right| \leq \sum\limits_{k = 0}^{\infty}{\left| x[ k] \right|  \left| e^{- j  \Omega  k} \right|} \leq \sum\limits_{k = 0}^{\infty}{X_{MAX}  r^{k}} = X_{MAX}  \sum\limits_{k = 0}^{\infty}r^{k} \label{eq:7.32}
\end{align}
Damit existiert die Fourier-Transformierte, wenn für den Betrag von r gilt:
\begin{align}
|r| < 1 \label{eq:7.33}
\end{align}
Also kann die Fourier-Transformierte für alle kausalen Signale, deren z-Transformierte \(X(z)\) ausschließlich Pole innerhalb des Einheitskreises besitzt, über die Definitionsgleichung berechnet werden.

\subsubsection{Zeitlich begrenzte Signalfolgen}\label{sec:zeitlich-begrenzte-signalfolgen}

Bei zeitlich begrenzten Signalfolgen mit endlicher Amplitude existiert immer eine Fourier-Transformierte. Zum Beweis wird die Definitionsgleichung der Fourier-Transformation nach oben abgeschätzt.
\begin{align}
\left| X(\Omega) \right| &= \left| \sum\limits_{k = - \infty}^{\infty}{x[ k]  e^{- j  \Omega  k}} \right| \leq \sum\limits_{k = - \infty}^{\infty}\left| x[ k] \right| = \sum\limits_{k = k_{1}}^{k_{2}}\left| x[ k] \right| \label{eq:7.34}
\end{align}
Da die Abtastwerte nur im Bereich \(k_{1}\) ≤ k ≤ \(k_{2}\) von null verschieden und ihre Werte begrenzt sind, ist die Summe endlich.

Bei der Berechnung der Fourier-Transformierten von Leistungssignalen in Abschnitt \ref{sec:7.1.6} zeigt sich, dass die Bedingung der absoluten Summierbarkeit eine hinreichende, aber keine notwendige Bedingung ist.

\subsection{Inverse Fourier-Transformation diskreter Folgen}\label{sec:inverse-fourier-transformation-diskreter-folgen}

Das Spektrum einer Signalfolge \(X(Ω)\) ist eine in 2π periodische Funktion. Ähnlich wie bei periodischen Zeitsignalen kann diese Funktion mit Hilfe einer Reihe dargestellt werden. Es ergibt sich die Folge der Abtastwerte
\begin{align}
x[ k] &= \frac{1}{2  \pi}  \int_{- \pi}^{\pi}{X(\Omega)  e^{j  \Omega  k}\ d\Omega} \label{eq:7.35}
\end{align}
Als Beweis wird gezeigt, dass die Fourier-Transformation von Folgen und ihre Rücktransformation invers zueinander sind. Dazu wird die Gleichung für die Fourier-Transformation von Folgen in die Gleichung zur Rücktransformation eingesetzt. Als Ergebnis muss sich die Folge \(x[k]\) ergeben.
\begin{align}
x[ k] &= \frac{1}{2  \pi}  \int_{- \pi}^{\pi}{X(\Omega)  e^{j  \Omega  k}\ d\Omega} = \frac{1}{2  \pi}  \int_{- \pi}^{\pi}{\left( \sum\limits_{\kappa = - \infty}^{\infty}{x[\kappa]  e^{- j  \Omega  \kappa}} \right)  e^{j  \Omega  k}\ d\Omega} \label{eq:7.36}
\end{align}
Zum Beweis wird die Reihenfolge von Summation und Integration vertauscht, und die beiden Exponentialfunktionen werden zusammengefasst.
\begin{align}
{x\left[k\right]} &=  {\frac{1}{2 \pi }  \int _{-\pi }^{\pi }\left(\sum\limits _{\kappa =-\infty }^{\infty }x\left[\kappa \right] e^{-j \Omega  \kappa }  \right) e^{j \Omega  k} {\; }d\Omega  =\sum\limits _{\kappa =-\infty }^{\infty }x\left[\kappa \right] \frac{1}{2 \pi }  \int _{-\pi }^{\pi }e^{-j \Omega  \kappa }  e^{j \Omega  k} {\; }d\Omega   } \\
{} &=  {\sum\limits _{\kappa =-\infty }^{\infty }x\left[\kappa \right] \frac{1}{2 \pi }  \int _{-\pi }^{\pi }e^{j \Omega  \left(k-\kappa \right)} {\; }d\Omega   }
 \label{eq:7.37}
\end{align}
Das Integral der komplexen Exponentialfunktion über eine volle Periode ist null. Nur für \(k = κ\) wird aus der komplexen Exponentialfunktion die Zahl eins. Das Ergebnis kann mit der Impulsfolge dargestellt werden. Mit der Ausblendeigenschaft der Impulsfolge ergibt sich
\begin{align}
\sum\limits_{\kappa = - \infty}^{\infty}{x[\kappa]  \frac{1}{2  \pi}  \int_{- \pi}^{\pi}{e^{j  \Omega  (k - \kappa)}\ d\Omega}} &= \sum\limits_{\kappa = - \infty}^{\infty}{x[\kappa]  \delta[ k - \kappa]} = x[ k] \label{eq:7.38}
\end{align}
Der Umgang mit der Rücktransformation wird an einem Beispiel vertieft.

\begin{example}{Inverse Fourier-Transformation diskreter Folgen}
\label{beispiel-inverse-fourier-transformation-diskreter-folgen}

Zu einem rechteckförmigen Spektrum der Form
\begin{align}
X(\Omega) &= \sigma\left( \Omega + \frac{\pi}{2} \right) - \sigma\left( \Omega - \frac{\pi}{2} \right) \label{eq:7.39}
\end{align}
soll die zugehörige Signalfolge \(x[k]\) berechnet werden. Dazu wird die Gleichung zur Rücktransformation verwendet. Für k ≠ 0 ergibt sich
\begin{align}
{x\left[k\right]} &=  {\frac{1}{2 \pi }  \int _{-\pi }^{\pi }X\left(\Omega \right) e^{j \Omega  k} {\; }d\Omega  =\frac{1}{2 \pi }  \int _{-\pi /2}^{\pi /2}e^{j \Omega  k} {\; }d\Omega  =\frac{1}{2 \pi }  \left. \frac{1}{j k}  e^{j \Omega  k} \right|_{-\pi /2}^{\pi /2} } \\
{} &=  {\frac{1}{2 \pi }  \frac{1}{j k}  2 j \sin \left(\frac{\pi }{2}  k\right)=\frac{1}{k \pi }  \sin \left(\frac{\pi }{2}  k\right)}
 \label{eq:7.40}
\end{align}
Für \(k = 0\) errechnet sich der Wert \(x[0]\) zu
\begin{align}
x[ 0] &= \frac{1}{2  \pi}  \int_{- \pi}^{\pi}{X\left( e^{j  \Omega} \right)  e^{j  \Omega  0}\ d\Omega} = \frac{1}{2  \pi}  \int_{- \pi/2}^{\pi/2}{1\ d\Omega} = \frac{\pi}{2  \pi} = \frac{1}{2} \label{eq:7.41}
\end{align}
Wie bei zeitkontinuierlichen Signalen kann aus dem Spektrum \(X(Ω)\) eine eindeutige Signalfolge \(x[k]\) berechnet werden. Bild \ref{fig:7.4} stellt die Signalfolge und ihre Fourier-Transformierte gegenüber.

\begin{fixedfigure}
  \includegraphics[alt={Signalfolge und ihr rechteckförmiges Spektrum}]{new media/7_4_signalfolge_und_ihr_rechteckfrmi.png}
  \captionof{figure}{Signalfolge und ihr rechteckförmiges Spektrum}
  \label{fig:7.4}
\end{fixedfigure}

\end{example}

\subsection{Fourier-Transformierte von Leistungssignalfolgen}\label{sec:7.1.6}

In Abschnitt \ref{sec:7.1.4} wird eine hinreichende Bedingung für die Existenz der Fourier-Transformierten einer Signalfolge hergeleitet und diskutiert. Es zeigt sich, dass die absolute Summierbarkeit eine hinreichende Bedingung ist. Aus diesem Grund existiert für Energiesignalfolgen immer eine Fourier-Transformierte. Aber auch für Leistungssignalfolgen lässt sich eine Fourier-Transformierte bestimmen. Ähnlich wie im zeitkontinuierlichen Bereich wird dazu die Impulsfunktion im Frequenzbereich verwendet. Da bei zeitdiskreten Signalen das Spektrum periodisch in 2π ist, wird die Signalfolge \(x[k]\) zu dem Spektrum
\begin{align}
X(\Omega) &= \sum\limits_{\nu = - \infty}^{\infty}{\delta(\Omega - 2  \pi  \nu)} \label{eq:7.42}
\end{align}
berechnet. Mit der Definitionsgleichung der inversen Fourier-Transformation von Signalfolgen ergibt sich
\begin{align}
x[ k] &= \frac{1}{2  \pi}  \int_{- \pi}^{\pi}{\sum\limits_{\nu = - \infty}^{\infty}{\delta(\Omega - 2  \pi  \nu)}  e^{j  \Omega  k}\ d\Omega} \label{eq:7.43}
\end{align}
In dem Bereich von - π ≤ Ω ≤ π liegt nur der Impuls an der Stelle \(Ω = 0\). Damit ergibt sich mit der Ausblendeigenschaft der Impulsfunktion
\begin{align}
x[ k] &= \frac{1}{2  \pi}  \int_{- \pi}^{\pi}{\delta(\Omega)  e^{j  \Omega  k}\ d\Omega} = \frac{1}{2  \pi} \label{eq:7.44}
\end{align}
Diese Folge \(x[k]\) ist nach den Ausführungen in Abschnitt \ref{sec:3.1.3} keine Energiesignalfolge, trotzdem besitzt sie eine Fourier-Transformierte. Bild \ref{fig:7.5} stellt die Einsfolge und ihre Fourier-Transformierte gegenüber.

\begin{figure}
  \centering
  \includegraphics[alt={Einsfolge und ihr Spektrum}]{new media/7_5_einsfolge_und_ihr_spektrum.png}
  \caption{Einsfolge und ihr Spektrum}
  \label{fig:7.5}
\end{figure}

Mit den Rechenregeln zur Fourier-Transformation von Signalfolgen lassen sich mit dieser Korrespondenz weitere Korrespondenzen herleiten.

\begin{example}{Korrespondenz der Kosinusfolge}
\label{beispiel-korrespondenz-der-kosinusfolge}

Bei der zeitkontinuierlichen Fourier-Transformation kann das Spektrum einer harmonischen Funktion über die Rücktransformation bestimmt werden. In diesem Beispiel wird diese Idee für Folgen aufgegriffen. Es wird die zu dem Spektrum
\begin{align}
X(\Omega) &= \sum\limits_{\nu = - \infty}^{\infty}{\delta\left( \Omega + \Omega_{0} + 2  \pi  \nu \right) + \delta\left( \Omega - \Omega_{0} + 2  \pi  \nu \right)} \label{eq:7.45}
\end{align}
gehörige Signalfolge \(x[k]\) berechnet. Dabei liegt \(Ω_{0}\) in dem Bereich 0 \textless{} \(Ω_{0}\) \textless{} π. Für die Berechnung der entsprechenden Signalfolge \(x[k]\) wird die Gleichung zur Rücktransformation verwendet und der Ausdruck in zwei separate Integrale aufgeteilt.
\begin{align}
{x\left[k\right]} &=  {\frac{1}{2 \pi }  \int _{-\pi }^{\pi }X\left(\Omega \right) e^{j \Omega  k} {\; }d\Omega  =\frac{1}{2 \pi }  \int _{-\pi }^{\pi }\left(\delta \left(\Omega +\Omega _{0} \right)+\delta \left(\Omega -\Omega _{0} \right)\right) e^{j \Omega  k} {\; }d\Omega  } \\
{} &=  {\frac{1}{2 \pi }  \int _{-\pi }^{\pi }\delta \left(\Omega +\Omega _{0} \right) e^{j \Omega  k} {\; }d\Omega  +\frac{1}{2 \pi }  \int _{-\pi }^{\pi }\delta \left(\Omega -\Omega _{0} \right) e^{j \Omega  k} {\; }d\Omega  }
 \label{eq:7.46}
\end{align}
Jedes Integral kann mit der Ausblendeigenschaft der Impulsfunktion gelöst werden.
\begin{align}
{x\left[k\right]} &=  {\frac{1}{2 \pi }  e^{-j \Omega _{0}  k}  \int _{-\pi }^{\pi }\delta \left(\Omega +\Omega _{0} \right){\; }d\Omega  +\frac{1}{2 \pi }  e^{j \Omega _{0}  k}  \int _{-\pi }^{\pi }\delta \left(\Omega -\Omega _{0} \right){\; }d\Omega  } \\
{} &=  {\frac{1}{2 \pi }  \left(e^{j \Omega _{0}  k} +e^{-j \Omega _{0}  k} \right)=\frac{2}{2 \pi }  \cos \left(\Omega _{0}  k\right)=\frac{1}{\pi }  \cos \left(k \Omega _{0} \right)}
 \label{eq:7.47}
\end{align}
Damit ergibt sich die Korrespondenz zur Kosinusfolge zu
\begin{align}
F\left\{ \cos\left( k  \Omega_{0} \right) \right\} &= \pi  \sum\limits_{\nu = - \infty}^{\infty}{\delta\left( \Omega + \Omega_{0} + 2  \pi  \nu \right) + \delta\left( \Omega - \Omega_{0} + 2  \pi  \nu \right)} \label{eq:7.48}
\end{align}
Bild \ref{fig:7.6} stellt eine Kosinusfolge mit \(Ω_{0} = π\)/3 und ihre Fourier-Transformierten dar.

\begin{fixedfigure}
  \includegraphics[alt={Kosinusfolge und ihr Spektrum}]{new media/7_6_kosinusfolge_und_ihr_spektrum.png}
  \captionof{figure}{Kosinusfolge und ihr Spektrum}
  \label{fig:7.6}
\end{fixedfigure}

Auf demselben Weg kann die Fourier-Transformierte der Sinusfolge hergeleitet werden.

\end{example}

\subsection{\texorpdfstring{Symmetrieeigenschaften der Fourier-Transformation von Signalfolgen }{Symmetrieeigenschaften der Fourier-Transformation von Signalfolgen }}\label{symmetrieeigenschaften-der-fourier-transformation-von-signalfolgen}

Bei der Anwendung der Fourier-Transformation können Symmetriebedingungen genutzt werden, um aus existierenden Korrespondenzen weitere Korrespondenzen abzuleiten. Insbesondere bei reellen Zahlenfolgen können die Symmetriebedingungen für eine schnelle Plausibilisierung von Ergebnissen verwendet werden.

Im Folgenden werden die wesentlichen Symmetrieeigenschaften für komplexe und reelle Signalfolgen zusammengefasst. Weitere Symmetrieregeln und ihre Herleitung sind in {[}Oppe04{]} zu finden. Dabei wird davon ausgegangen, dass zur Signalfolge \(x[k]\) das Spektrum \(X(Ω)\) gehört.

\begin{table}
  \centering
  \caption{Symmetrieregeln der Fourier-Transformation für komplexe Signalfolgen}
  \begin{tabular}{ll}
    \toprule
    Folge \(x[k]\) & Fourier-Transformierte \(X(Ω)\) \\
    \midrule
    Gerade Folge \(x[k]\) & Reelles Spektrum \(X(Ω)\) \\
    Ungerade Folge \(x[k]\) & Imaginäres Spektrum \(X(Ω)\) \\
    Konjugiert komplexe Folge x*{[}k{]} & X*(- Ω) \\
    Gespiegelte konjugiert komplexe Folge x*{[}- k{]} & X*(Ω) \\
    \bottomrule
  \end{tabular}
\end{table}

\begin{table}
  \centering
  \caption{Symmetrieregeln der Fourier-Transformation für reelle Signalfolgen \(x[k]\)}
  \label{tab:7.3}
  \begin{tabular}{ll}
    \toprule
    Regel & Anmerkung \\
    \midrule
    \(X(Ω)\) = X*(- Ω) & 
Fourier-Transformierte einer reellen Folge 
ist konjugiert symmetrisch \\
    \(X_{R}(Ω)\) = \(X_{R}\)(- Ω) & 
Realteil der Fourier-Transformierten 
ist gerade \\
    \(X_{I}(Ω)\) = - \(X_{I}\)(- Ω) & 
Imaginärteil der Fourier-Transformierten 
ist ungerade \\
    \textbar \(X(Ω)\)\textbar{} = \textbar X(- Ω)\textbar{} & 
Betrag der Fourier-Transformierten 
ist gerade \\
    ϕ(Ω) = - ϕ(- Ω) & 
Phase der Fourier-Transformierten 
ist ungerade \\
    \bottomrule
  \end{tabular}
\end{table}

Die hier dargestellten Symmetrieregeln für reelle Signalfolgen werden an einem Beispiel verdeutlicht.

\begin{example}{Fourier-Transformierten der kausalen Exponentialfolge}
\label{beispiel-fourier-transformierten-der-kausalen-exponentialfolge}

Es wird die Fourier-Transformierte der kausalen Exponentialfolge
\begin{align}
x[ k] &= \lambda^{k}  \sigma[ k] \label{eq:7.49}
\end{align}
diskutiert. Das Spektrum der kausalen Exponentialfolge berechnet sich zu
\begin{align}
X(\Omega) &= \frac{1}{1 - \lambda  e^{- j  \Omega}} = \frac{1}{1 - \lambda  \cos(\Omega) - j  \lambda  \sin(\Omega)} \label{eq:7.50}
\end{align}
Die konjugiert symmetrische Fourier-Transformierte ergibt sich zu
\begin{align}
X^{*}( - \Omega) &= \frac{1}{1 - \lambda  \cos( - \Omega) + j  \lambda  \sin( - \Omega)} = \frac{1}{1 - \lambda  \cos(\Omega) - j  \lambda  \sin(\Omega)} \label{eq:7.51}
\end{align}
Beide Transformierten stimmen erwartungsgemäß überein. Die Spektren der beiden Folgen sind in Bild \ref{fig:7.7} sowohl als Real- und Imaginärteil, als auch als Betrag und Phase dargestellt.

\begin{fixedfigure}
  \includegraphics[alt={Spektren der kausalen Exponentialfolgen mit \(λ = 0.5\) und \(λ = 0.9\)}]{new media/7_7_spektren_der_kausalen_exponentia.png}
  \captionof{figure}{Spektren der kausalen Exponentialfolgen mit \(λ = 0.5\) und \(λ = 0.9\)}
  \label{fig:7.7}
\end{fixedfigure}

Der Realteil des Spektrums ist gerade, der Imaginärteil ist ungerade. Der Betrag des Spektrums ist gerade, die Phase ungerade. Damit entsprechen die Symmetrien den in Tabelle \ref{tab:7.3} dargestellten Symmetrieregeln.

\end{example}

\section{Rechenregeln der Fourier-Transformation von Signalfolgen}\label{sec:7.2}

Wie andere Transformationen besitzt die Fourier-Transformation von Signalfolgen Rechenregeln, die die Berechnung weiterer Korrespondenzen vereinfachen. Die Regeln werden im Folgenden hergeleitet und jeweils an einem Beispiel verdeutlicht. Sie entsprechen weitgehend den Rechenregeln für die Fourier-Transformation zeitkontinuierlicher Signale. Am Ende dieses Abschnittes sind alle Rechenregeln in einer Übersicht zusammengestellt.

\subsection{Linearität}\label{sec:linearituxe4t-3}

Wie die übrigen Integraltransformationen ist auch die Fourier-Transformation von Signalfolgen eine lineare Transformation. Durch Einsetzen in die Definitionsgleichung ergibt sich
\begin{align}
F\left\{ \nu_{1}  x_{1}[ k] + \nu_{2}  x_{2}[ k] \right\} &= \sum\limits_{k = - \infty}^{\infty}{\left( \nu_{2}  x_{1}[ k] + \nu_{2}  x_{2}[ k] \right)  e^{- j  \Omega  k}} \label{eq:7.52}
\end{align}
Im Fall absolut summierbarer Reihen besitzen beide Folgen \(x_{1}\){[}k{]} und \(x_{2}\){[}k{]} eigene Fourier-Transformierte \(X_{1}(Ω)\) und \(X_{2}(Ω)\). In dem Fall gilt das Distributivgesetz
\begin{align}
{\sum\limits_{k = - \infty}^{\infty}{\left( \nu_{1}  x_{1}[ k] + \nu_{2}  x_{2}[ k] \right)  e^{- j  k  \Omega}} = \sum\limits_{k = - \infty}^{\infty}{\nu_{1}  x_{1}[ k]  e^{- j  \Omega  k}} + \sum\limits_{k = - \infty}^{\infty}{\nu_{2}  x_{2}[ k]  e^{- j  \Omega  k}}
}{			 = \nu_{1}  X_{1}(\Omega) + \nu_{2}  X_{2}(\Omega)} \label{eq:7.53}
\end{align}
Damit ist die Linearitätseigenschaft für Signalfolgen, die eine Fourier-Transformierte besitzen, bewiesen.

\begin{example}{Linearität}
\label{beispiel-linearituxe4t-1}

Mit Hilfe der Linearität kann zum Beispiel die Fourier-Transformierte der Summe zweier kausaler Exponentialfolgen errechnet werden.
\begin{align}
F\left\{ 3  \frac{1}{2^{k}}  \sigma[ k] - 2  \frac{1}{5^{k}}  \sigma[ k] \right\} &= \frac{3}{1 - \frac{1}{2}  e^{- j  \Omega}} - \frac{2}{1 - \frac{1}{5}  e^{- j  \Omega}} = \frac{6}{2 - 1  e^{- j  \Omega}} - \frac{10}{5 - 1  e^{- j  \Omega}} \label{eq:7.54}
\end{align}
\end{example}

\subsection{Verschiebung im Zeitbereich}\label{sec:verschiebung-im-zeitbereich}

Ist das Spektrum \(X(Ω)\) einer Signalfolge \(x[k]\) bekannt, ergibt sich das Spektrum der verschobenen Folge durch Einsetzen in die Definitionsgleichung zu
\begin{align}
F\left\{ x\left[ k - k_{0} \right] \right\} &= \sum\limits_{k = - \infty}^{\infty}{x\left[ k - k_{0} \right]  e^{- j  \Omega  k}} \label{eq:7.55}
\end{align}
Da die Summation von - ∞ bis + ∞ durchgeführt wird, kann der Index k um \(k_{0}\) verschoben werden. Es ergibt sich
\begin{align}
{F\left\{ x\left[ k - k_{0} \right] \right\} = \sum\limits_{k = - \infty}^{\infty}{x\left[ k - k_{0} \right]  e^{- j  \Omega  k}} = \sum\limits_{k = - \infty}^{\infty}{x[ k]  e^{- j  \Omega  \left( k + k_{0} \right)}}
}{	 = e^{- j  \Omega  k_{0}}  \sum\limits_{k = - \infty}^{\infty}{x[ k]  e^{- j  \Omega  k}} = e^{- j  \Omega  k_{0}}  X(\Omega)} \label{eq:7.56}
\end{align}
Das Spektrum ändert nur seine Phase, der Betrag bleibt unverändert.

\begin{example}{Verschiebung im Zeitbereich kombiniert mit Linearität}
\label{beispiel-verschiebung-im-zeitbereich-kombiniert-mit-linearituxe4t}

Als Beispiel wird die Fourier-Transformierte der in Bild \ref{fig:7.8} dargestellten Signalfolge berechnet. Sie berechnet sich aus der Differenz zweier Rechteckfolgen der Länge \(K = 5\), von denen die erste um\\
\(k_{1} = 1\) nach rechts und die zweite Rechteckfolge um \(k_{2}\) = - 5 nach links verschoben ist.
\begin{align}
x[ k] &= \sigma[ k - 1] - \sigma[ k - 6] - \left( \sigma[ k + 5] - \sigma[ k] \right) \label{eq:7.57}
\end{align}
Entsprechend ergibt sich die Fourier-Transformierte zu
\begin{align}
{X\left(\Omega \right)} &=  {e^{-j \frac{4 \Omega }{2} }  \frac{\sin \left(\frac{5 \Omega }{2} \right)}{\sin \left(\frac{\Omega }{2} \right)}  e^{-j \Omega } -e^{-j \frac{4 \Omega }{2} }  \frac{\sin \left(\frac{5 \Omega }{2} \right)}{\sin \left(\frac{\Omega }{2} \right)}  e^{-j \left(-5\right) \Omega } } \\
{} &=  {\frac{\sin \left(\frac{5 \Omega }{2} \right)}{\sin \left(\frac{\Omega }{2} \right)}  \left(e^{-j 3 \Omega } -e^{j 3 \Omega } \right)}
 \label{eq:7.58}
\end{align}
Der Betrag des Spektrums ist in Bild \ref{fig:7.8} dargestellt.

\begin{fixedfigure}
  \includegraphics[alt={Signalfolge und Betrag des Spektrums für das Beispiel\\
Verschiebung im Zeitbereich kombiniert mit Linearität}]{new media/7_8_signalfolge_und_betrag_des_spekt.png}
  \captionof{figure}{Signalfolge und Betrag des Spektrums für das Beispiel\\
Verschiebung im Zeitbereich kombiniert mit Linearität}
  \label{fig:7.8}
\end{fixedfigure}

\end{example}

\subsection{Verschiebung im Frequenzbereich}\label{sec:verschiebung-im-frequenzbereich}

Eine Verschiebung kann nicht nur im Zeitbereich, sondern auch im Frequenzbereich erfolgen. In diesem Fall ergibt sich die zugehörige Zeitfunktion aus
\begin{align}
{X\left( \Omega - \Omega_{0} \right) = \sum\limits_{k = - \infty}^{\infty}{x[ k]  e^{- j  k  \left( \Omega - \Omega_{0} \right)}} = \sum\limits_{k = - \infty}^{\infty}{x[ k]  e^{- j  k  \Omega}  e^{j  k  \Omega_{0}}}
}{	 = \sum\limits_{k = - \infty}^{\infty}{x[ k]  e^{j  k  \Omega_{0}}  e^{- j  k  \Omega}} = F\left\{ x[ k]  e^{j  k  \Omega_{0}} \right\}} \label{eq:7.59}
\end{align}
\begin{example}{Frequenzverschiebung}
\label{beispiel-frequenzverschiebung}

Typisches Beispiel für eine Frequenzverschiebung ist eine Modulation. Dabei wird eine Folge \(x[k]\) mit einer harmonischen Folge multipliziert.
\begin{align}
y[ k] &= x[ k]  \cos\left( \Omega_{0}  k \right) = x[ k]  \frac{1}{2}  \left( e^{j  \Omega_{0}  k} + e^{- j  \Omega_{0}  k} \right) \label{eq:7.60}
\end{align}
Durch die Multiplikation verschiebt sich das ursprüngliche Spektrum \(X(Ω)\) um \(Ω_{0}\) nach links und rechts.
\begin{align}
Y(\Omega) &= \frac{1}{2}  \left( X\left( \Omega + \Omega_{0} \right) + X\left( \Omega - \Omega_{0} \right) \right) \label{eq:7.61}
\end{align}
Das Spektrum wird damit aus dem ursprünglichen Spektralbereich in einen Spektralbereich um ± \(Ω_{0}\) verschoben. Bild \ref{fig:7.9} stellt die Auswirkung einer Modulation auf den Spektralbereich grafisch dar.

\begin{fixedfigure}
  \includegraphics[alt={Beispiel für die Verschiebung eines Spektralbereiches durch Modulation}]{new media/7_9_beispiel_fr_die_verschiebung_ein.png}
  \captionof{figure}{Beispiel für die Verschiebung eines Spektralbereiches durch Modulation}
  \label{fig:7.9}
\end{fixedfigure}

\end{example}

\subsection{Zeitliche Spiegelung}\label{sec:zeitliche-spiegelung}

Ist das Spektrum \(X(Ω)\) einer Signalfolge \(x[k]\) bekannt, ergibt sich das Spektrum der Folge \(x[- k]\) durch Einsetzen in die Definitionsgleichung zu
\begin{align}
F\left\{ x[ - k] \right\} &= \sum\limits_{k = - \infty}^{\infty}{x[ - k]  e^{- j  \Omega  k}} = \sum\limits_{n = - \infty}^{\infty}{x[ n]  e^{j  \Omega  n}} = X( - \Omega) \label{eq:7.62}
\end{align}
\begin{example}{Zeitliche Spiegelung}
\label{beispiel-zeitliche-spiegelung}

Als Anwendungen der zeitlichen Spiegelung wird das Spektrum der Folge
\begin{align}
x[ k] &= \lambda^{|k|} \label{eq:7.63}
\end{align}
für \textbar λ\textbar{} \textless{} 1 berechnet werden. Das Spektrum der einseitigen Exponentialfolge wird über die Definitionsgleichung berechnet.
\begin{align}
F\left\{ \lambda^{k}  \sigma[ k] \right\} &= \sum\limits_{k = 0}^{\infty}\left( \lambda  e^{- j  \Omega} \right)^{k} = \frac{1}{1 - \lambda  e^{- j  \Omega}}\ \label{eq:7.64}
\end{align}
Diese Folge wird achsensymmetrisch ergänzt. Das Signal ergibt sich durch die Summe einer kausalen Exponentialfunktion, der gespiegelten Funktion und einer Impulsfolge, die den Wert an der Stelle \(k = 0\) korrigiert.
\begin{align}
x[ k] &= \lambda^{k}  \sigma[ k] + \lambda^{- k}  \sigma[ - k] - \delta[ k] \label{eq:7.65}
\end{align}
Das Spektrum errechnet sich mit der Regel der zeitlichen Spiegelung und der Linearität zu
\begin{align}
X(\Omega) &= \frac{1}{1 - \lambda  e^{- j  \Omega}} + \frac{1}{1 - \lambda  e^{j  \Omega}} - 1 \label{eq:7.66}
\end{align}
Die Folge \(x[k]\) und der Betrag des Spektrums sind in Bild \ref{fig:7.10} dargestellt.

\begin{fixedfigure}
  \includegraphics[alt={Achsensymmetrische Exponentialfolge und ihr Spektrum für \(λ = 0.7\)}]{new media/7_10_achsensymmetrische_exponentialf.png}
  \captionof{figure}{Achsensymmetrische Exponentialfolge und ihr Spektrum für \(λ = 0.7\)}
  \label{fig:7.10}
\end{fixedfigure}

\end{example}
\subsection{Differenz von Folgen}\label{sec:differenz-von-folgen}

Ist das Spektrum \(X(Ω)\) einer Signalfolge \(x[k]\) bekannt, ergibt sich das Spektrum der Differenz zweier aufeinanderfolgender Folgenwerte \(x[k]\) - \(x[k - 1]\) durch Anwendung der Verschiebungsregel zu
\begin{align}
F\left\{ x[ k] - x[ k - 1] \right\} &= X(\Omega)  \left( 1 - e^{- j  \Omega} \right) \label{eq:7.67}
\end{align}
Die Differenzregel ist damit eine Anwendung der Verschiebungsregel. Da die Differenz aber der Differentiation zeitkontinuierlicher Signale entspricht, wird sie als eigene Rechenregel dargestellt.

\subsection{Summe von Folgen}\label{sec:summe-von-folgen}

Ist das Spektrum \(X(Ω)\) einer Signalfolge \(x[k]\) bekannt, ergibt sich das Spektrum der Summe aller Folgenwerte zu
\begin{align}
F\left\{ \sum\limits_{\kappa = - \infty}^{k}{x[\kappa]} \right\} &= \frac{1}{1 - e^{- j  \Omega}}  X(\Omega) + \pi  X(0)  \sum\limits_{\nu = - \infty}^{\infty}{\delta(\Omega + 2  \pi  \nu)} \label{eq:7.68}
\end{align}
Diese Regel wird ohne Beweis und Beispiel aufgeführt. Der Beweis wird über die Faltungsregel in Kombination mit der Fourier-Transformierten einer Sprungfolge in einer Übungsaufgabe geführt.

\subsection{Differentiation im Frequenzbereich}\label{sec:differentiation-im-frequenzbereich}

Wird das Spektrum \(X(Ω)\) einer Signalfolge \(x[k]\) abgeleitet, wird die Folge im Zeitbereich linear gewichtet. Es ergibt sich die Rechenregel
\begin{align}
F\left\{ k  x[ k] \right\} &= j  \frac{dX}{d\Omega} \label{eq:7.69}
\end{align}
Der Beweis ergibt sich aus einer Umformung der Definitionsgleichung.
\begin{align}
\frac{dX}{d\Omega} &= \frac{d}{d\Omega}\ \sum\limits_{k = - \infty}^{\infty}{x[ k]  e^{- j  \Omega  k}} = \sum\limits_{k = - \infty}^{\infty}{- j  k  x[ k]  e^{- j  \Omega  k}} = - j  \sum\limits_{k = - \infty}^{\infty}{k  x[ k]  e^{- j  \Omega  k}} \label{eq:7.70}
\end{align}
\begin{example}{Differentiation im Frequenzbereich}
\label{beispiel-differentiation-im-frequenzbereich}

Die Rechenregel zur Differentiation im Frequenzbereich kann dazu verwendet werden, das Spektrum der Folge
\begin{align}
x[ k] &= k  \left( \sigma[ k + K] - \sigma\left[ k - (K + 1) \right] \right) = k  y[ k] \label{eq:7.71}
\end{align}
zu berechnen. Mit der Regel für die Differentiation im Frequenzbereich und dem Spektrum der Rechteckfolge \(y[k]\) ergibt sich
\begin{align}
X(\Omega) &= - j  \frac{dY}{d\Omega} = - j  \frac{d}{d\Omega}\frac{\sin(K  \Omega)}{\sin\left( \frac{\Omega}{2} \right)} = - j  \frac{K  \cos(K  \Omega)  \sin\left( \frac{\Omega}{2} \right) - \frac{1}{2}  \cos\left( \frac{\Omega}{2} \right)  \sin(K  \Omega)}{\sin^{2}\left( \frac{\Omega}{2} \right)} \label{eq:7.72}
\end{align}
\end{example}

\subsection{Faltung im Zeitbereich}\label{sec:faltung-im-zeitbereich}

Das Ausgangssignal zeitdiskreter LTI-Systeme errechnet sich über die Faltungssumme.
\begin{align}
y[ k] &= g[ k]*u[ k] = \sum\limits_{\kappa = - \infty}^{\infty}{g[\kappa]  u[ k - \kappa]} = \sum\limits_{\kappa = - \infty}^{\infty}{g[ k - \kappa]}  u[\kappa] \label{eq:7.73}
\end{align}
Sind die Spektren \(G(Ω)\) der Signalfolge \(g[k]\) und \(U(Ω)\) der Signalfolge \(u[k]\) bekannt, errechnet sich das Spektrum des Ausgangssignals über das Produkt der beiden Spektren
\begin{align}
Y(\Omega) &= G(\Omega)  U(\Omega) \label{eq:7.74}
\end{align}
Der Beweis kann wie bei der z-Transformation durch Einsetzen in die Definitionsgleichung geführt werden.
\begin{align}
F\left\{ g[ k]*u[ k] \right\} &= F\left\{ \sum\limits_{\kappa = - \infty}^{\infty}{g[\kappa]  u[ k - \kappa]} \right\} = \sum\limits_{k = - \infty}^{\infty}{\sum\limits_{\kappa = - \infty}^{\infty}{g[\kappa]  u[ k - \kappa]}}  e^{- j  \Omega  k} \label{eq:7.75}
\end{align}
Vertauschen der Summationsreihenfolge und Substitution ergibt
\begin{align}
{\boldsymbol{\mathit{F}}\left\{g\left[k\right]*u\left[k\right]\right\}} &=  {\sum\limits _{k=-\infty }^{\infty }\sum\limits _{\kappa =-\infty }^{\infty }g\left[\kappa \right] u\left[k-\kappa \right]   e^{-j \Omega  k} =\sum\limits _{\kappa =-\infty }^{\infty }\sum\limits _{k=-\infty }^{\infty }g\left[\kappa \right] u\left[k-\kappa \right]   e^{-j \Omega  k} } \\
{} &=  {\sum\limits _{\kappa =-\infty }^{\infty }g\left[\kappa \right] \sum\limits _{k=-\infty }^{\infty }u\left[k-\kappa \right]   e^{-j \Omega  k}  e^{j \Omega  \kappa }  e^{-j \Omega  \kappa } } \\
{} &=  {\sum\limits _{n=-\infty }^{\infty }g\left[\kappa \right] e^{-j \Omega  \kappa }  \sum\limits _{k=-\infty }^{\infty }u\left[k-\kappa \right]   e^{-j \Omega  \left(k-\kappa \right)} =\sum\limits _{n=-\infty }^{\infty }g\left[n\right] e^{-j \Omega  n}  \sum\limits _{m=-\infty }^{\infty }u\left[m\right]   e^{-j \Omega  m} } \\
{} &=  {G\left(\Omega \right) U\left(\Omega \right)}
 \label{eq:7.76}
\end{align}
Die Faltungsregel ist eine wesentliche Rechenregel der Fourier-Transformation für Signalfolgen, da sich das Ausgangssignal eines linearen zeitinvarianten Systems aus der Faltung von Eingangsfolge und Impulsantwort ergibt. Die Faltungsregel wird im Kapitel \hyperref[section]{1} detailliert aufgegriffen.

\subsection{Multiplikation im Zeitbereich}\label{sec:7.2.9}

Neben der Faltung ist bei der digitalen Signalverarbeitung die Multiplikation im Zeitbereich eine wichtige Operation. Sie wird verwendet, um den Ausschnitt eines Signals zu beschreiben, indem das Signal \(x[k]\) mit einer sogenannten Fensterfolge \(w[k]\) multipliziert wird. Die Rechenregel zur Multiplikation im Zeitbereich ermöglicht die Berechnung des Spektrums des Signalausschnittes.

Sind die Spektren \(X(Ω)\) einer Signalfolge \(x[k]\) und \(W(Ω)\) einer Signalfolge \(w[k]\) bekannt, errechnet sich das Spektrum des Signals, das aus dem Produkt der beiden Signale
\begin{align}
y[ k]\  &= \ x[ k]  w[ k] \label{eq:7.77}
\end{align}
entsteht, durch Einsetzen der Signale in die Definitionsgleichung
\begin{align}
Y(\Omega) &= \sum\limits_{k = - \infty}^{\infty}{x[ k]  w[ k]}  e^{- j  \Omega  k} \label{eq:7.78}
\end{align}
Die Substitution der Folge \(w[k]\) durch ihre inverse Fourier-Transformierte
\begin{align}
Y(\Omega) &= \sum\limits_{k = - \infty}^{\infty}{x[ k]  \frac{1}{2  \pi}  \int_{- \pi}^{\pi}{W(\Theta)  e^{j  k  \Theta}\ d\Theta}}  e^{- j  \Omega  k} \label{eq:7.79}
\end{align}
und Tauschen der Reihenfolge von Integration und Summation
\begin{align}
Y(\Omega) &= \sum\limits_{k = - \infty}^{\infty}{x[ k]  \frac{1}{2  \pi}  \int_{- \pi}^{\pi}{W(\Theta)  e^{j  k  \Theta}\ d\Theta}}  e^{- j  \Omega  k} = \frac{1}{2  \pi}  \int_{- \pi}^{\pi}{W(\Theta)  \sum\limits_{k = - \infty}^{\infty}{x[ k]  e^{j  (\Theta - \Omega)  k}\ }d\Theta} \label{eq:7.80}
\end{align}
führt unter Anwendung der Verschiebungsregel im Frequenzbereich zu
\begin{align}
{Y\left(\Omega \right)} &=  {\frac{1}{2 \pi }  \int _{-\pi }^{\pi }W\left(\Theta \right) \sum\limits _{k=-\infty }^{\infty }x\left[k\right] e^{j \left(\Theta -\Omega \right) k} {\; } d\Theta  =\frac{1}{2 \pi }  \int _{-\pi }^{\pi }W\left(\Theta \right) X\left(\Omega -\Theta \right){\; }d\Theta  } \\
{} &=  {\frac{1}{2 \pi }  X\left(\Omega \right)*W\left(\Omega \right)}
 \label{eq:7.81}
\end{align}
Die Multiplikation zweier Folgen im Zeitbereich führt demnach zu einer Faltung der entsprechenden Spektren im Frequenzbereich.

\begin{example}{Fensterung als Multiplikation im Zeitbereich}
\label{beispiel-fensterung-als-multiplikation-im-zeitbereich}

Als Beispiel wird eine Kosinusfolge betrachtet, von der ein Ausschnitt über zwei Periodendauern vorliegt. Der Signalausschnitt ist in Bild \ref{fig:7.11} dargestellt und kann mathematisch beschrieben werden als
\begin{align}
{y[ k]\  = \ w[ k]  x[ k] = \left( \sigma[ k + 10] - \sigma[ k - 11] \right)  \cos\left( \frac{2  \pi}{10}  k \right)
}{= \left( \sigma[ k + 10] - \sigma[ k - 11] \right)  \cos\left( \frac{\pi}{5}  k \right)} \label{eq:7.82}
\end{align}
Das Spektrum der Folge \(w[k]\) ergibt sich aus der berechneten Rechteckfolge und der Verschiebungsregel zu
\begin{align}
W(\Omega) &= \frac{\sin\left( \frac{21  \Omega}{2} \right)}{\sin\left( \frac{\Omega}{2} \right)} \label{eq:7.83}
\end{align}
Das Spektrum der Kosinusfolge \(x[k]\) wird berechnet zu
\begin{align}
X(\Omega) &= \pi  \left( \delta\left( \Omega + \frac{\pi}{5} \right) + \delta\left( \Omega - \frac{\pi}{5} \right) \right) \label{eq:7.84}
\end{align}
Die Multiplikation im Zeitbereich führt zu Faltung im Frequenzbereich. Die Faltung der beiden Spektren führt zu einer Verschiebung von \(W(Ω)\) an die Stellen der Impulse. Damit ergibt sich für das Spektrum \(Y(Ω)\)
\begin{align}
{Y\left(\Omega \right)} &=  {\frac{1}{2 \pi }  \pi  \left(\frac{\sin \left(\frac{21 \left(\Omega +\frac{\pi }{5} \right)}{2} \right)}{\sin \left(\frac{\left(\Omega +\frac{\pi }{5} \right)}{2} \right)} +\frac{\sin \left(\frac{21 \left(\Omega -\frac{\pi }{5} \right)}{2} \right)}{\sin \left(\frac{\left(\Omega -\frac{\pi }{5} \right)}{2} \right)} \right)} \\
{} &=  {\frac{1}{2}  \left(\frac{\sin \left(21 \left(\frac{\Omega }{2} +\frac{\pi }{10} \right)\right)}{\sin \left(\frac{\Omega }{2} +\frac{\pi }{10} \right)} +\frac{\sin \left(21 \left(\frac{\Omega }{2} -\frac{\pi }{10} \right)\right)}{\sin \left(\frac{\Omega }{2} -\frac{\pi }{10} \right)} \right)}
 \label{eq:7.85}
\end{align}
Signal und Betrag des Spektrums sind in Bild \ref{fig:7.11} dargestellt.

\begin{fixedfigure}
  \includegraphics[alt={Ausschnitt einer Kosinusfolge und das entsprechende Spektrum}]{new media/7_11_ausschnitt_einer_kosinusfolge_u.png}
  \captionof{figure}{Ausschnitt einer Kosinusfolge und das entsprechende Spektrum}
  \label{fig:7.11}
\end{fixedfigure}

Durch die Multiplikation der Folge mit einer Rechteckfolge wird das Signal ausgeschnitten. Im Frequenzbereich falten sich die Spektren der beiden Signale, wobei das Spektrum der Kosinusfunktion aus zwei Impulsen an den Stellen ± \(Ω_{0}\) besteht. Durch die Faltung wird das Spektrum der Rechteckfolge an die Stellen ± \(Ω_{0}\) verschoben. Der Effekt, der bei der Fensterung von Signalen entsteht, wird bei der Schätzung von Spektren in Kapitel \hyperref[diskrete-fourier-transformation]{10} weiter vertieft. Zusätzlich wird das Spektrum wegen der Abtastung periodisch in 2π wiederholt.

\end{example}

\subsection{Parsevalsche Gleichung}\label{sec:parsevalsche-gleichung}

Die Parsevalsche Gleichung drückt aus, dass die Summe über dem Quadrat einer Folge gleich dem Integral über dem Quadrat ihrer Transformierten ist. Wie bei zeitkontinuierlichen Signalen wird es dazu verwendet, die Energie eines Signals im Zeitbereich oder im Frequenzbereich zu berechnen.
\begin{align}
E &= \sum\limits_{k = - \infty}^{\infty}{x[ k]  x^{*}[ k]} = \frac{1}{2  \pi}  \int_{- \pi}^{\pi}{X(\Omega)  X^{*}(\Omega)\ d\Omega} \label{eq:7.86}
\end{align}
Der Beweis dieser Rechenregel ergibt sich aus der Faltungsregel. Der linke Teil der Gleichung entspricht der Faltungssumme von \(x[k]\) und x*{[}- k{]} an der Stelle \(k = 0\).
\begin{align}
\left. \ x[ k]*x^{*}[ - k] \right|_{k = 0} &= \left. \ \sum\limits_{\kappa = - \infty}^{\infty}{x[\kappa]  x^{*}\left[ - (k - \kappa) \right]} \right|_{k = 0} = \left. \ \sum\limits_{\kappa = - \infty}^{\infty}{x[\kappa]  x^{*}[ - k + \kappa]} \right|_{k = 0} = \sum\limits_{\kappa = - \infty}^{\infty}{x[\kappa]  x^{*}[\kappa]} \label{eq:7.87}
\end{align}
Der Faltung im Zeitbereich entspricht die Multiplikation im Frequenzbereich. Mit den Spektren
\begin{align}
x[ k] \origof X(\Omega) \label{eq:7.88}
\end{align}
und
\begin{align}
x^{*}[ - k] \origof X^{*}(\Omega) \label{eq:7.89}
\end{align}
sowie der Definitionsgleichung für die inverse Fourier-Transformierte an der Stelle \(k = 0\) ergibt sich
\begin{align}
\sum\limits_{\kappa = - \infty}^{\infty}{x[\kappa]  x^{*}[\kappa]} &= \left. \ \frac{1}{2  \pi}  \int_{- \pi}^{\pi}{X(\Omega)  X^{*}(\Omega)  e^{j  k  \Omega}\ d\Omega} \right|_{k = 0} = \frac{1}{2  \pi}  \int_{- \pi}^{\pi}{X(\Omega)  X^{*}(\Omega)\ d\Omega} \label{eq:7.90}
\end{align}
Der Ausdruck X(Ω)X*(Ω) entspricht dem Betragsquadrat \textbar \(X(Ω)\)\textbar{}\textsuperscript{2} und wird als Energiedichtespektrum bezeichnet. Es ist insbesondere bei der Berechnung stochastischer Prozesse von Bedeutung und wird in Teil C dieser Buchreihe wieder aufgegriffen.

\subsection{Zusammenfassung der Rechenregeln zur Fourier-Transformation von Signalfolgen}\label{sec:zusammenfassung-der-rechenregeln-zur-fourier-transformation-von-signalfolgen}

Zur besseren Übersicht stellt Tabelle \ref{tab:7.4} die wesentlichen Eigenschaften der Fourier-Transformation von Signalfolgen noch einmal zusammen.

\begin{table}
  \centering
  \caption{Rechenregeln der Fourier-Transformation von Signalfolgen}
  \label{tab:7.4}
  \begin{tabular}{lll}
    \toprule
    \textbf{Regel} & \textbf{Folge \(x[k]\)} & \textbf{Fourier-Transformierte \(X(Ω)\)} \\
    \midrule
    Linearität & \(\nu_{1}  x_{1}\lbrack k\rbrack + \nu_{2}  x_{2}\lbrack k\rbrack\) & \(\nu_{1}  X_{1}(\Omega) + \nu_{2}  X_{2}(\Omega)\) \\
    Zeitverschiebung & \(x\left\lbrack k - k_{0} \right\rbrack\) & \(e^{- j  \Omega  k_{0}}  X(\Omega)\) \\
    Frequenzverschiebung & \(x\lbrack k\rbrack  e^{j  k  \Omega_{0}}\) & \(X\left( \Omega - \Omega_{0} \right)\) \\
    Zeitliche Spiegelung & \(x\lbrack - k\rbrack\) & \(X( - \Omega)\) \\
    Differenz & \(x\lbrack k\rbrack - x\lbrack k - 1\rbrack\) & \(X(\Omega)  \left( 1 - e^{- j  \Omega} \right)\) \\
    Summation & \(\sum\limits_{\kappa = - \infty}^{k}{x\lbrack\kappa\rbrack}\) & \(\frac{1}{1 - e^{- j  \Omega}}  X(\Omega) + \pi  X(0)  \sum\limits_{\nu = - \infty}^{\infty}{\delta(\Omega + 2  \pi  \nu)}\) \\
    lineare Gewichtung & \(k  x\lbrack k\rbrack\) & \(j  \frac{dX}{d\Omega}\) \\
    Faltung & \(g\lbrack k\rbrack*u\lbrack k\rbrack\) & \(G(\Omega)  U(\Omega)\) \\
    Multiplikation & \(\ x\lbrack k\rbrack  w\lbrack k\rbrack\) & \(\frac{1}{2  \pi}  X(\Omega)*W(\Omega)\) \\
    Parsevalsche Gleichung & \(E = \sum\limits_{k = - \infty}^{\infty}{x\lbrack k\rbrack  x^{*}\lbrack k\rbrack}\) & \(E = \frac{1}{2  \pi}  \int_{- \pi}^{\pi}{X(\Omega)  X^{*}(\Omega)\ d\Omega}\) \\
    \bottomrule
  \end{tabular}
\end{table}

Mit Hilfe der zusammengestellten Rechenregeln lassen sich weitere Korrespondenzen bestimmen.

\begin{example}{Fourier-Transformierte der Sprungfolge}
\label{beispiel-fourier-transformierte-der-sprungfolge}

Die Rechenregeln der zeitdiskreten Fourier-Transformation werden angewendet, um die Korrespondenz der Sprungfolge zu bestimmen. Die Sprungfolge kann in zwei Anteile zerlegt werden, in die konstante Folge
\begin{align}
x_{1}[ k] &= 0.5 \label{eq:7.91}
\end{align}
und die Folge
\begin{align}
x_{2}[ k] = \left\{ \begin{aligned}
 & - 0.5	für\ k < 0 \\
 & + 0.5	für\ k \geq 0
\end{aligned} \right.\ \label{eq:7.92}
\end{align}
\begin{fixedfigure}
  \includegraphics[alt={Zerlegung der Sprungfolge in eine konstante Folge und eine Signum-Folge}]{new media/7_12_zerlegung_der_sprungfolge_in_ei.png}
  \captionof{figure}{Zerlegung der Sprungfolge in eine konstante Folge und eine Signum-Folge}
\end{fixedfigure}

Die Summe der beiden Folgen ergibt die Sprungfolge. Für den Teil \(x_{1}\){[}k{]} kann die Fourier-Transformierte mit der Korrespondenz der Einsfolge und der Linearität angegeben werden zu
\begin{align}
X_{1}(\Omega) &= \pi  \sum\limits_{\nu = - \infty}^{\infty}{\delta(\Omega - 2  \pi  \nu)} \label{eq:7.93}
\end{align}
Zur Ermittlung der Fourier-Transformierten des zweiten Teils \(x_{2}\){[}k{]} wird die Impulsfolge verwendet. Sie kann dargestellt werden als
\begin{align}
\delta[ k] &= x_{2}[ k] - x_{2}[ k - 1] \label{eq:7.94}
\end{align}
Die Fourier-Transformierte der Impulsfolge ist 1 und eine Verschiebung der Funktion um \(k_{0} = 1\) führt zur Multiplikation der Fourier-Transformierten mit \(e^{-jΩ}\). Damit ergibt sich für die Fourier-Transformierte
\begin{align}
1 &= \left( 1 - e^{- j  \Omega} \right)  \sum\limits_{k = - \infty}^{\infty}{x_{2}[ k]  e^{- j  \Omega  k}} = \left( 1 - e^{- j  \Omega} \right)  X_{2}\left( e^{j  \Omega} \right) \label{eq:7.95}
\end{align}
und durch Auflösen
\begin{align}
X_{2}(\Omega) &= \frac{1}{1 - e^{- j  \Omega}} \label{eq:7.96}
\end{align}
Durch Superposition von Gleichung \eqref{eq:7.93} und Gleichung \eqref{eq:7.96} ergibt sich die Fourier-Transformierte des diskreten Einheitssprungs zu
\begin{align}
F\left( \sigma[ k] \right) &= \frac{1}{1 - e^{- j  \Omega}} + \pi  \sum\limits_{\nu = - \infty}^{\infty}{\delta(\Omega - 2  \pi  \nu)} \label{eq:7.97}
\end{align}
Bild \ref{fig:7.13} stellt den diskreten Einheitssprung und den Betrag seiner Fourier-Transformierten gegenüber.

\begin{fixedfigure}
  \includegraphics[alt={Diskreter Einheitssprung und der Betrag seines Spektrums}]{new media/7_13_diskreter_einheitssprung_und_de.png}
  \captionof{figure}{Diskreter Einheitssprung und der Betrag seines Spektrums}
  \label{fig:7.13}
\end{fixedfigure}

\end{example}

\subsection{Korrespondenzen der Fourier-Transformation von Signalfolgen}\label{sec:korrespondenzen-der-fourier-transformation-von-signalfolgen}

Tabelle \ref{tab:7.5} stellt wichtige Korrespondenzen der Fourier-Transformation von Signalfolgen zusammen.

\begin{table}
  \centering
  \caption{Korrespondenzen der Fourier-Transformation von Signalfolgen}
  \label{tab:7.5}
  \begin{tabular}{lll}
    \toprule
    Nr. & Folge \(x[k]\) & Fourier-Transformierte \(X(Ω)\) \\
    \midrule
    1 & \(\delta\lbrack k\rbrack\) & \(1\) \\
    2 & \(1\) & \(2  \pi  \sum\limits_{\nu = - \infty}^{\infty}{\delta(\Omega + 2  \pi  \nu)}\) \\
    3 & \(\delta\left\lbrack k - k_{0} \right\rbrack\) & \(e^{- j  \Omega  k_{0}}\) \\
    4 & \(\lambda^{k}  \sigma\lbrack k\rbrack\) & \(\frac{1}{1 - \lambda  e^{- j  \Omega}}\text{        für  }|\lambda| < 1\) \\
    5 & \((k + 1)  \lambda^{k}  \sigma\lbrack k\rbrack\) & \(\frac{1}{\left( 1 - \lambda  e^{- j  \Omega} \right)^{2}}\text{   für  }|\lambda| < 1\) \\
    6 & \(\sigma\lbrack k\rbrack\) & \(\frac{1}{1 - e^{- j  \Omega}} + \pi  \sum\limits_{\nu = - \infty}^{\infty}{\delta(\Omega + 2  \pi  \nu)}\) \\
    7 & \(e^{j  \Omega  k_{0}}\) & \(2  \pi  \sum\limits_{\nu = - \infty}^{\infty}{\delta\left( \Omega - \Omega_{0} + 2  \pi  \nu \right)}\) \\
    8 & \(x\lbrack k\rbrack = \sigma\lbrack k\rbrack - \sigma\left\lbrack k - k_{0} \right\rbrack\) & \(e^{- j  \frac{\Omega  \left( k_{0} - 1 \right)}{2}}  \frac{\sin\left( \frac{k_{0}  \Omega}{2} \right)}{\sin\left( \frac{\Omega}{2} \right)}\) \\
    9 & \(\sin\left\lbrack \Omega_{0}  k \right\rbrack\) & \(j  \pi  \sum\limits_{n = - \infty}^{\infty}{\delta\left( \Omega + \Omega_{0} + 2  \pi  n \right) - \delta\left( \Omega - \Omega_{0} + 2  \pi  n \right)}\) \\
    10 & \(\cos\left\lbrack \Omega_{0}  k \right\rbrack\) & \(\pi  \sum\limits_{n = - \infty}^{\infty}{\delta\left( \Omega + \Omega_{0} + 2  \pi  n \right) + \delta\left( \Omega - \Omega_{0} + 2  \pi  n \right)}\) \\
    11 & \({sgn}\lbrack k\rbrack\) & \(\frac{2}{1 - e^{- j  \Omega}}\) \\
    \bottomrule
  \end{tabular}
\end{table}

\section{Fourier-Transformation von Signalfolgen und andere Integraltrans\-formationen}\label{sec:7.3}

Bei der Diskussion der Fourier-Transformation von Signalfolgen in den vorangegangenen Abschnitten wird auf die Ähnlichkeit zur Fourier-Transformation von Signalfolgen und der z-Transformation verwiesen. Zur Verbesserung des Verständnisses wird der Zusammenhang zwischen der Fourier-Transformation von Signalfolgen und den anderen Integraltransformationen weiter vertieft.

\subsection{Zusammenhang zwischen Fourier-Transformation von Signalfolgen und Fourier-Transformation kontinuierlicher Signale}\label{sec:zusammenhang-zwischen-fourier-transformation-von-signalfolgen-und-fourier-transformation-kontinuierlicher-signale}

In Kapitel \ref{sec:7.1.2} wird die Definitionsgleichung der Fourier-Transformation für Folgen über die Fourier-Transformation von ideal abgetasteten Signalen hergeleitet. Der Zusammenhang zwischen den unterschiedlichen Signalen und ihren Spektren wird in diesem Abschnitt weiter vertieft. In Bild \ref{fig:7.14} sind die unterschiedlichen Signale und ihre Spektren dargestellt.

\begin{figure}
  \centering
  \includesvg{new media/7_14_zusammenhang_zwischen_den_spekt.svg}
  \caption{Zusammenhang zwischen den Spektren des zeitkontinuierlichen Signals \(x(t)\),\\
des abgetasteten Signals \(x_{A}(t)\) und der Fourier-Transformierten der entsprechenden Folge \(x_{D}\){[}k{]}}
  \label{fig:7.14}
\end{figure}

Das zeitkontinuierliche Signal \(x(t)\) besitzt ein Spektrum \(X(ω)\). Um das Abtasttheorem zu erfüllen, wird es mit einem Tiefpass bandbegrenzt. Es entsteht das Signal \(x_{TP}(t)\) mit dem Spektrum \(X_{TP}(ω)\). Wegen der Abtastung ist das Spektrum \(X_{A}(ω)\) periodisch in \(ω_{A}\) und mit dem Kehrwert der Abtastzeit 1/\(T_{A}\) skaliert. Die Folge \(x_{D}\){[}k{]} der Abtastwerte besitzt das Spektrum \(X_{D}(Ω)\), das für Ω = ωT\textsubscript{A} dem Spektrum \(X_{A}(ω)\) entspricht. Das Basisband von dem Spektrum \(X_{A}(ω)\) entspricht dem Spektrum \(X_{TP}(ω)\) und, wenn das Signal \(x(t)\) bandbegrenzt ist, auch dem Spektrum \(X(ω)\). Dieser Zusammenhang der verschiedenen Signale und ihrer Spektren wird an einem Beispiel illustriert.

\begin{example}{Zusammenhang zwischen Spektren zeitkontinuierlicher Signale und Signalfolgen}
\label{beispiel-zusammenhang-zwischen-spektren-zeitkontinuierlicher-signale-und-signalfolgen}

Gegeben ist ein kausales, zeitkontinuierliches Signal
\begin{align}
x(t) &= \left( e^{- t} - e^{- 2  t} \right)  \sigma(t) \label{eq:7.98}
\end{align}
Das Signal besitzt die Laplace-Transformierte
\begin{align}
X(s) &= \frac{1}{s + 1} - \frac{1}{s + 2} = \frac{1}{s + 1}  \frac{1}{s + 2} \label{eq:7.99}
\end{align}
Da die Laplace-Transformierte nur Pole in der negativen Halbebene besitzt, ergibt sich das Spektrum \(X(ω)\) zu
\begin{align}
X(\omega) &= \frac{1}{j  \omega + 1}  \frac{1}{j  \omega + 2} = \frac{1}{2 - \omega^{2} + 3  j  \omega} \label{eq:7.100}
\end{align}
Das zeitkontinuierliche Signal \(x(t)\) sowie der Betrag des zugehörigen Spektrums \textbar \(X(ω)\)\textbar{} sind in Bild \ref{fig:7.15} dargestellt.

\begin{fixedfigure}
  \includegraphics[alt={Zeitkontinuierliches Signal \(x(t)\) sowie der Betrag des zugehörigen Spektrums \textbar \(X(ω)\)\textbar{}}]{new media/7_15_zeitkontinuierliches_signal_xt.png}
  \captionof{figure}{Zeitkontinuierliches Signal \(x(t)\) sowie der Betrag des zugehörigen Spektrums \textbar \(X(ω)\)\textbar{}}
  \label{fig:7.15}
\end{fixedfigure}

Das Signal wird mit einer Abtastzeit \(T_{A} = π\)/25 abgetastet. Dadurch wird das Spektrum periodisch in \(ω_{A} = 50\) wiederholt und mit dem Faktor 1/\(T_{A} = 25\)/π multipliziert. Das Spektrum des ideal abgetasteten Signals \(x_{A}(t)\) ist in Bild \ref{fig:7.16} dargestellt.

\begin{fixedfigure}
  \includegraphics[alt={Ideal abgetastetes Signal \(x_{A}(t)\) sowie der Betrag des zugehörigen Spektrums \textbar \(X_{A}(ω)\)\textbar{}}]{new media/7_16_ideal_abgetastetes_signal_xat_s.png}
  \captionof{figure}{Ideal abgetastetes Signal \(x_{A}(t)\) sowie der Betrag des zugehörigen Spektrums \textbar \(X_{A}(ω)\)\textbar{}}
  \label{fig:7.16}
\end{fixedfigure}

Das Spektrum der zugehörigen Signalfolge
\begin{align}
x_{D}[ k] &= \left( e^{- k  T_{A}} - e^{- 2  k  T_{A}} \right)  \sigma[ k] \label{eq:7.101}
\end{align}
ist kontinuierlich und errechnet sich über die Definitionsgleichung der Fourier-Transformation für Folgen zu
\begin{align}
X_{D}(\Omega) &= \sum\limits_{k = - \infty}^{\infty}{x_{D}[ k]  e^{- j  \Omega  k}} = \sum\limits_{k = - \infty}^{\infty}{\left( e^{- k  T_{A}} - e^{- 2  k  T_{A}} \right)  \sigma[ k]  e^{- j  \Omega  k}} \label{eq:7.102}
\end{align}
Da es sich um eine kausale Folge handelt, kann die Summe auf zwei geometrische Reihen zurückgeführt und berechnet werden.
\begin{align}
{X_{D} \left(\Omega \right)} &=  {\sum\limits _{k=0}^{\infty }\left(e^{-k T_{A} } -e^{-2 k T_{A} } \right) e^{-j \Omega  k}  =\sum\limits _{k=0}^{\infty }e^{-k T_{A} }  e^{-j \Omega  k}  -\sum\limits _{k=0}^{\infty }e^{-2 k T_{A} }  e^{-j \Omega  k}  } \\
{} &=  {\sum\limits _{k=0}^{\infty }\left(e^{-T_{A} -j \Omega } \right)^{k}  -\sum\limits _{k=0}^{\infty }\left(e^{-2 T_{A} -j \Omega } \right)^{k}  =\frac{1}{1-e^{-T_{A} -j \Omega } } -\frac{1}{1-e^{-2 T_{A} -j \Omega } } }
 \label{eq:7.103}
\end{align}
Bild \ref{fig:7.17} zeigt die Signalfolge \(x_{D}\){[}k{]} und den Betrag des Spektrums der Signalfolge \textbar \(X_{D}(Ω)\)\textbar{} im Vergleich zum Spektrum \(X_{A}(ω)\) des ideal abgetasteten Signals \(x_{A}(t)\), wobei \(X_{A}(ω)\) mit normierter Frequenz dargestellt ist.

\begin{fixedfigure}
  \includegraphics[alt={Folge der Abtastwert \(x_{D}\){[}k{]} sowie der Betrag des zugehörigen Spektrums \textbar \(X_{D}(Ω)\)\textbar{}}]{new media/7_17_folge_der_abtastwert_xdk_sowie.png}
  \captionof{figure}{Folge der Abtastwert \(x_{D}\){[}k{]} sowie der Betrag des zugehörigen Spektrums \textbar \(X_{D}(Ω)\)\textbar{}}
  \label{fig:7.17}
\end{fixedfigure}

Ein Vergleich der unterschiedlichen Spektren zeigt, dass mit der Fourier-Transformation der Signalfolge \(x_{D}\){[}k{]} das Spektrum des zeitkontinuierlichen Signals berechnet werden kann, wenn das Abtasttheorem eingehalten wird. Allerdings müssen dazu unendlich viele Abtastwerte \(x_{D}\){[}k{]} vorliegen, was normalerweise nicht der Fall ist. Die Bestimmung des Spektrums eines Signals bei einer endlichen Anzahl von Abtastwerten wird in Kapitel \hyperref[diskrete-fourier-transformation]{10} diskutiert.

\end{example}

\subsection{Zusammenhang zwischen z-Transformation und Fourier-Transformation kausaler Signalfolgen}\label{sec:7.3.2}

Bei dem Vergleich der Definitionsgleichungen der z-Transformation
\begin{align}
X(z) &= \sum\limits_{k = 0}^{\infty}{x[ k]  z^{- k}} \label{eq:7.104}
\end{align}
und der Fourier-Transformation von kausalen Folgen
\begin{align}
X(\Omega) &= \sum\limits_{k = 0}^{\infty}{x[ k]  e^{- j  k  \Omega}} \label{eq:7.105}
\end{align}
fällt auf, dass die Variable z der z-Transformation bei der Fourier-Transformation von kausalen Folgen durch den Ausdruck \(e^{jΩ}\) ersetzt wird. Grafisch gesehen wird bei dem Übergang von der\\
z-Transformierten auf die Fourier-Transformierte das Spektrum über dem Einheitskreis abgewickelt.

\begin{example}{Zusammenhang zwischen z-Transformation und Fourier-Transformation von Signalfolgen}
\label{beispiel-zusammenhang-zwischen-z-transformation-und-fourier-transformation-von-signalfolgen}

Dieser Sachverhalt wird anhand eines Signals mit der z-Transformierten
\begin{align}
X(z) &= \frac{z^{2}}{z^{2} - z + 0.5} \label{eq:7.106}
\end{align}
anschaulich dargestellt. Bild \ref{fig:7.18} zeigt die Pol-Nullstellenverteilung der z-Transformierten \(X(z)\).

\begin{fixedfigure}
  \includegraphics[alt={Pol-Nullstellenverteilung der z-Transformierten \(X(z)\)}]{new media/7_18_polnullstellenverteilung_der_zt.png}
  \captionof{figure}{Pol-Nullstellenverteilung der z-Transformierten \(X(z)\)}
  \label{fig:7.18}
\end{fixedfigure}

Die z-Transformierte \(X(z)\) hat ein konjugiertes Polpaar an der Stelle \(α = 0.5\) ± 0.5j und eine doppelte Nullstelle \(β = 0\). Beide Pole liegen innerhalb des Einheitskreises. Der Betrag der z-Transformierten und der Betrag der Fourier-Transformierten werden in Bild \ref{fig:7.19} verglichen.

\begin{fixedfigure}
  \includegraphics[alt={Vergleich des Betrags der z-Transformierten \(X(z)\) und der Fourier-Transformierten}]{new media/7_19_vergleich_des_betrags_der_ztran.png}
  \captionof{figure}{Vergleich des Betrags der z-Transformierten \(X(z)\) und der Fourier-Transformierten}
  \label{fig:7.19}
\end{fixedfigure}

In der linken Grafik ist der Betrag der z-Transformierten dargestellt. An den beiden Polen wird der Betrag der Übertragungsfunktion unendlich groß, in der Grafik wird er auf den Wert 4 begrenzt. Außerdem ist mit dem schwarzen Kreis der Betrag der Übertragungsfunktion über dem Einheitskreis dargestellt. Die Abwicklung dieses Kreises entspricht dem in der rechten Grafik dargestellten Betrag des Spektrums. Der Betrag des Spektrums ist nahe den Polen groß und sinkt mit steigendem Abstand zu den Polen ab.

Bild \ref{fig:7.20} stellt die Phase der z-Transformierten und die Phase der Fourier-Transformierten dar.

\begin{fixedfigure}
  \includegraphics[alt={Vergleich der Phase der z-Transformierten \(X(z)\) und der Fourier-Transformierten}]{new media/7_20_vergleich_der_phase_der_ztransf.png}
  \captionof{figure}{Vergleich der Phase der z-Transformierten \(X(z)\) und der Fourier-Transformierten}
  \label{fig:7.20}
\end{fixedfigure}

In der linken Grafik ist die Phase der z-Transformierten dargestellt. Zwischen der Nullstelle und den Polen im Koordinatenursprung ändert sich die Phase von - π auf π. Aufgrund der räumlichen Darstellung ist dieser Sprung leider verdeckt. Mit dem schwarzen Kreis ist die Phase der Übertragungsfunktion über dem Einheitskreis dargestellt. Die Abwicklung dieses Kreises entspricht der in der rechten Grafik dargestellten Phase des Spektrums. Die Phasenänderung des Spektrums ist nahe den Polen groß und sinkt mit steigendem Abstand zu den Polen ab.

\end{example}

Der Vergleich zeigt, dass die Fourier-Transformierte einer Folge aus der z-Transformierten über die Substitution z = \(e^{jΩ}\) gewonnen werden kann. Dazu ist es allerdings erforderlich, dass der Einheitskreis im Konvergenzbereich der z-Transformierten liegt. Nach den Ausführungen zur z-Transformation in Kapitel \ref{sec:5.1} ist das der Fall, wenn alle Pole der z-Transformierten im Einheitskreis liegen.

\begin{example}{Berechnung der Fourier-Transformierten aus der z-Transformierten}
\label{beispiel-berechnung-der-fourier-transformierten-aus-der-z-transformierten}

Die Fourier-Transformierte eines Signals mit der z-Transformierten
\begin{align}
X(z) &= \frac{z^{2}}{z^{2} - z + 0.5} \label{eq:7.107}
\end{align}
soll berechnet werden. Die Pole der z-Transformierten liegen nach Bild \ref{fig:7.18} innerhalb des Einheitskreises. Damit ergibt sich die Fourier-Transformierte aus
\begin{align}
X(\Omega) &= \left. \ X(z) \right|_{z = e^{j  \Omega}} = \frac{e^{j  2  \Omega}}{e^{j  2  \Omega} - e^{j  \Omega} + 0.5} \label{eq:7.108}
\end{align}
Der Betrag des Spektrums ist bereits in Bild \ref{fig:7.19} dargestellt.

\end{example}

\section{Berechnung von Korrespondenzen der Fourier-Transformation von Signalfolgen}\label{sec:7.4}

In den vorangegangenen Abschnitten werden unterschiedliche Methoden zur Berechnung von Korrespondenzen der Fourier-Transformation von Signalfolgen diskutiert. Zur besseren Übersicht stellt Bild \ref{fig:7.21} die unterschiedlichen Methoden zusammen und gibt die Bedingungen an, unter denen die entsprechenden Methoden angewendet werden können.

\begin{figure}
  \centering
  \includesvg{new media/7_21_methoden_zur_berechnung_von_kor.svg}
  \caption{Methoden zur Berechnung von Korrespondenzen der Fourier-Transformation von Signalfolgen}
  \label{fig:7.21}
\end{figure}

Korrespondenzen der Fourier-Transformation können über die Definitionsgleichungen der Fourier-Transformation und der inversen Fourier-Transformation berechnet werden, wenn die entsprechenden Integrale konvergieren. Bei Energiesignalfolgen ist das für die Definitionsgleichung der Fourier-Transformation generell der Fall. Bei Leistungssignalfolgen wird ausgehend von einer Fourier-Transformierten \(X(Ω)\) die zugehörige Signalfolge \(x[k]\) bestimmt.

Aus der z-Transformierten \(X(z)\) ergibt sich die Fourier-Transformierte \(X(Ω)\) durch Substitution, wenn die Signalfolge \(x[k]\) kausal ist und die Pole von \(X(z)\) im Einheitskreis (\textbar α\textbar{} \textless{} 1).
